%!TeX root=../princesstop.tex
\chapter{Becky}

\begin{figure}[t!]
\centering
\includegraphics[width=\linewidth]{filigree}
\end{figure}

\lettrine[lines=5]{O}{f} course the greatest power Sara possessed and the one which gained her even more followers than her luxuries and the fact that she was »the show pupil,« the power that Lavinia and certain other girls were most envious of, and at the same time most fascinated by in spite of themselves, was her power of telling stories and of making everything she talked about seem like a story, whether it was one or not.

Anyone who has been at school with a teller of stories knows what the wonder means—how he or she is followed about and besought in a whisper to relate romances; how groups gather round and hang on the outskirts of the favoured party in the hope of being allowed to join in and listen. Sara not only could tell stories, but she adored telling them. When she sat or stood in the midst of a circle and began to invent wonderful things, her green eyes grew big and shining, her cheeks flushed, and, without knowing that she was doing it, she began to act and made what she told lovely or alarming by the raising or dropping of her voice, the bend and sway of her slim body, and the dramatic movement of her hands. She forgot that she was talking to listening children; she saw and lived with the fairy folk, or the kings and queens and beautiful ladies, whose adventures she was narrating. Sometimes when she had finished her story, she was quite out of breath with excitement, and would lay her hand on her thin, little, quick-rising chest, and half laugh as if at herself.

»When I am telling it,« she would say, »it doesn't seem as if it was only made up. It seems more real than you are—more real than the schoolroom. I feel as if I were all the people in the story—one after the other. It is queer.«

She had been at Miss Minchin's school about two years when, one foggy winter's afternoon, as she was getting out of her carriage, comfortably wrapped up in her warmest velvets and furs and looking very much grander than she knew, she caught sight, as she crossed the pavement, of a dingy little figure standing on the area steps, and stretching its neck so that its wide-open eyes might peer at her through the railings. Something in the eagerness and timidity of the smudgy face made her look at it, and when she looked she smiled because it was her way to smile at people.

But the owner of the smudgy face and the wide-open eyes evidently was afraid that she ought not to have been caught looking at pupils of importance. She dodged out of sight like a jack-in-the-box and scurried back into the kitchen, disappearing so suddenly that if she had not been such a poor little forlorn thing, Sara would have laughed in spite of herself. That very evening, as Sara was sitting in the midst of a group of listeners in a corner of the schoolroom telling one of her stories, the very same figure timidly entered the room, carrying a coal box much too heavy for her, and knelt down upon the hearth rug to replenish the fire and sweep up the ashes.

She was cleaner than she had been when she peeped through the area railings, but she looked just as frightened. She was evidently afraid to look at the children or seem to be listening. She put on pieces of coal cautiously with her fingers so that she might make no disturbing noise, and she swept about the fire irons very softly. But Sara saw in two minutes that she was deeply interested in what was going on, and that she was doing her work slowly in the hope of catching a word here and there. And realizing this, she raised her voice and spoke more clearly.

»The Mermaids swam softly about in the crystal-green water, and dragged after them a fishing-net woven of deep-sea pearls,« she said. »The Princess sat on the white rock and watched them.«

It was a wonderful story about a princess who was loved by a Prince Merman, and went to live with him in shining caves under the sea.

The small drudge before the grate swept the hearth once and then swept it again. Having done it twice, she did it three times; and, as she was doing it the third time, the sound of the story so lured her to listen that she fell under the spell and actually forgot that she had no right to listen at all, and also forgot everything else. She sat down upon her heels as she knelt on the hearth rug, and the brush hung idly in her fingers. The voice of the storyteller went on and drew her with it into winding grottos under the sea, glowing with soft, clear blue light, and paved with pure golden sands. Strange sea flowers and grasses waved about her, and far away faint singing and music echoed.

The hearth brush fell from the work-roughened hand, and Lavinia Herbert looked round.

»That girl has been listening,« she said.

The culprit snatched up her brush, and scrambled to her feet. She caught at the coal box and simply scuttled out of the room like a frightened rabbit.

Sara felt rather hot-tempered.

»I knew she was listening,« she said. »Why shouldn't she?«

Lavinia tossed her head with great elegance.

»Well,« she remarked, »I do not know whether your mamma would like you to tell stories to servant girls, but I know \textsc{my} mamma wouldn't like \textsc{me} to do it.«

»My mamma!« said Sara, looking odd. »I don't believe she would mind in the least. She knows that stories belong to everybody.«

»I thought,« retorted Lavinia, in severe recollection, »that your mamma was dead. How can she know things?«

»Do you think she \textsc{doesn't} know things?« said Sara, in her stern little voice. Sometimes she had a rather stern little voice.

»Sara's mamma knows everything,« piped in Lottie. »So does my mamma—'cept Sara is my mamma at Miss Minchin's—my other one knows everything. The streets are shining, and there are fields and fields of lilies, and everybody gathers them. Sara tells me when she puts me to bed.«

»You wicked thing,« said Lavinia, turning on Sara; »making fairy stories about heaven.«

»There are much more splendid stories in Revelation,« returned Sara. »Just look and see! How do you know mine are fairy stories? But I can tell you«—with a fine bit of unheavenly temper—»you will never find out whether they are or not if you're not kinder to people than you are now. Come along, Lottie.« And she marched out of the room, rather hoping that she might see the little servant again somewhere, but she found no trace of her when she got into the hall.

»Who is that little girl who makes the fires?« she asked Mariette that night.

Mariette broke forth into a flow of description.

Ah, indeed, Mademoiselle Sara might well ask. She was a forlorn little thing who had just taken the place of scullery maid—though, as to being scullery maid, she was everything else besides. She blacked boots and grates, and carried heavy coal-scuttles up and down stairs, and scrubbed floors and cleaned windows, and was ordered about by everybody. She was fourteen years old, but was so stunted in growth that she looked about twelve. In truth, Mariette was sorry for her. She was so timid that if one chanced to speak to her it appeared as if her poor, frightened eyes would jump out of her head.

»What is her name?« asked Sara, who had sat by the table, with her chin on her hands, as she listened absorbedly to the recital.

Her name was Becky. Mariette heard everyone below-stairs calling, »Becky, do this,« and »Becky, do that,« every five minutes in the day.

Sara sat and looked into the fire, reflecting on Becky for some time after Mariette left her. She made up a story of which Becky was the ill-used heroine. She thought she looked as if she had never had quite enough to eat. Her very eyes were hungry. She hoped she should see her again, but though she caught sight of her carrying things up or down stairs on several occasions, she always seemed in such a hurry and so afraid of being seen that it was impossible to speak to her.

But a few weeks later, on another foggy afternoon, when she entered her sitting room she found herself confronting a rather pathetic picture. In her own special and pet easy-chair before the bright fire, Becky—with a coal smudge on her nose and several on her apron, with her poor little cap hanging half off her head, and an empty coal box on the floor near her—sat fast asleep, tired out beyond even the endurance of her hard-working young body. She had been sent up to put the bedrooms in order for the evening. There were a great many of them, and she had been running about all day. Sara's rooms she had saved until the last. They were not like the other rooms, which were plain and bare. Ordinary pupils were expected to be satisfied with mere necessaries. Sara's comfortable sitting room seemed a bower of luxury to the scullery maid, though it was, in fact, merely a nice, bright little room. But there were pictures and books in it, and curious things from India; there was a sofa and the low, soft chair; Emily sat in a chair of her own, with the air of a presiding goddess, and there was always a glowing fire and a polished grate. Becky saved it until the end of her afternoon's work, because it rested her to go into it, and she always hoped to snatch a few minutes to sit down in the soft chair and look about her, and think about the wonderful good fortune of the child who owned such surroundings and who went out on the cold days in beautiful hats and coats one tried to catch a glimpse of through the area railing.

On this afternoon, when she had sat down, the sensation of relief to her short, aching legs had been so wonderful and delightful that it had seemed to soothe her whole body, and the glow of warmth and comfort from the fire had crept over her like a spell, until, as she looked at the red coals, a tired, slow smile stole over her smudged face, her head nodded forward without her being aware of it, her eyes drooped, and she fell fast asleep. She had really been only about ten minutes in the room when Sara entered, but she was in as deep a sleep as if she had been, like the Sleeping Beauty, slumbering for a hundred years. But she did not look—poor Becky—like a Sleeping Beauty at all. She looked only like an ugly, stunted, worn-out little scullery drudge.

Sara seemed as much unlike her as if she were a creature from another world.

On this particular afternoon she had been taking her dancing lesson, and the afternoon on which the dancing master appeared was rather a grand occasion at the seminary, though it occurred every week. The pupils were attired in their prettiest frocks, and as Sara danced particularly well, she was very much brought forward, and Mariette was requested to make her as diaphanous and fine as possible.

Today a frock the colour of a rose had been put on her, and Mariette had bought some real buds and made her a wreath to wear on her black locks. She had been learning a new, delightful dance in which she had been skimming and flying about the room, like a large rose-coloured butterfly, and the enjoyment and exercise had brought a brilliant, happy glow into her face.

When she entered the room, she floated in with a few of the butterfly steps—and there sat Becky, nodding her cap sideways off her head.

»Oh!« cried Sara, softly, when she saw her. »That poor thing!«

It did not occur to her to feel cross at finding her pet chair occupied by the small, dingy figure. To tell the truth, she was quite glad to find it there. When the ill-used heroine of her story wakened, she could talk to her. She crept toward her quietly, and stood looking at her. Becky gave a little snore.

»I wish she'd waken herself,« Sara said. »I don't like to waken her. But Miss Minchin would be cross if she found out. I'll just wait a few minutes.«

She took a seat on the edge of the table, and sat swinging her slim, rose-coloured legs, and wondering what it would be best to do. Miss Amelia might come in at any moment, and if she did, Becky would be sure to be scolded.

»But she is so tired,« she thought. »She is so tired!«

A piece of flaming coal ended her perplexity for her that very moment. It broke off from a large lump and fell on to the fender. Becky started, and opened her eyes with a frightened gasp. She did not know she had fallen asleep. She had only sat down for one moment and felt the beautiful glow—and here she found herself staring in wild alarm at the wonderful pupil, who sat perched quite near her, like a rose-coloured fairy, with interested eyes.

She sprang up and clutched at her cap. She felt it dangling over her ear, and tried wildly to put it straight. Oh, she had got herself into trouble now with a vengeance! To have impudently fallen asleep on such a young lady's chair! She would be turned out of doors without wages.

She made a sound like a big breathless sob.

»Oh, miss! Oh, miss!« she stuttered. »I arst yer pardon, miss! Oh, I do, miss!«

Sara jumped down, and came quite close to her.

»Don't be frightened,« she said, quite as if she had been speaking to a little girl like herself. »It doesn't matter the least bit.«

»I didn't go to do it, miss,« protested Becky. »It was the warm fire—an' me bein' so tired. It—it \textsc{wasn't} impertience!«

Sara broke into a friendly little laugh, and put her hand on her shoulder.

»You were tired,« she said; »you could not help it. You are not really awake yet.«

How poor Becky stared at her! In fact, she had never heard such a nice, friendly sound in anyone's voice before. She was used to being ordered about and scolded, and having her ears boxed. And this one—in her rose-coloured dancing afternoon splendour—was looking at her as if she were not a culprit at all—as if she had a right to be tired—even to fall asleep! The touch of the soft, slim little paw on her shoulder was the most amazing thing she had ever known.

»Ain't—ain't yer angry, miss?« she gasped. »Ain't yer goin' to tell the missus?«

»No,« cried out Sara. »Of course I'm not.«

The woeful fright in the coal-smutted face made her suddenly so sorry that she could scarcely bear it. One of her queer thoughts rushed into her mind. She put her hand against Becky's cheek.

»Why,« she said, »we are just the same—I am only a little girl like you. It's just an accident that I am not you, and you are not me!«

Becky did not understand in the least. Her mind could not grasp such amazing thoughts, and »an accident« meant to her a calamity in which some one was run over or fell off a ladder and was carried to »the 'orspital.«

»A' accident, miss,« she fluttered respectfully. »Is it?«

»Yes,« Sara answered, and she looked at her dreamily for a moment. But the next she spoke in a different tone. She realized that Becky did not know what she meant.

»Have you done your work?« she asked. »Dare you stay here a few minutes?«

Becky lost her breath again.

»Here, miss? Me?«

Sara ran to the door, opened it, and looked out and listened.

»No one is anywhere about,« she explained. »If your bedrooms are finished, perhaps you might stay a tiny while. I thought—perhaps—you might like a piece of cake.«

The next ten minutes seemed to Becky like a sort of delirium. Sara opened a cupboard, and gave her a thick slice of cake. She seemed to rejoice when it was devoured in hungry bites. She talked and asked questions, and laughed until Becky's fears actually began to calm themselves, and she once or twice gathered boldness enough to ask a question or so herself, daring as she felt it to be.

»Is that\longdash« she ventured, looking longingly at the rose-coloured frock. And she asked it almost in a whisper. »Is that there your best?«

»It is one of my dancing-frocks,« answered Sara. »I like it, don't you?«

For a few seconds Becky was almost speechless with admiration. Then she said in an awed voice, »Onct I see a princess. I was standin' in the street with the crowd outside Covin' Garden, watchin' the swells go inter the operer. An' there was one everyone stared at most. They ses to each other, »That's the princess.« She was a growed-up young lady, but she was pink all over—gownd an' cloak, an' flowers an' all. I called her to mind the minnit I see you, sittin' there on the table, miss. You looked like her.«

»I've often thought,« said Sara, in her reflecting voice, »that I should like to be a princess; I wonder what it feels like. I believe I will begin pretending I am one.«

Becky stared at her admiringly, and, as before, did not understand her in the least. She watched her with a sort of adoration. Very soon Sara left her reflections and turned to her with a new question.

»Becky,« she said, »weren't you listening to that story?«

»Yes, miss,« confessed Becky, a little alarmed again. »I knowed I hadn't orter, but it was that beautiful I—I couldn't help it.«

»I liked you to listen to it,« said Sara. »If you tell stories, you like nothing so much as to tell them to people who want to listen. I don't know why it is. Would you like to hear the rest?«

Becky lost her breath again.

»Me hear it?« she cried. »Like as if I was a pupil, miss! All about the Prince—and the little white Mer-babies swimming about laughing—with stars in their hair?«

Sara nodded.

»You haven't time to hear it now, I'm afraid,« she said; »but if you will tell me just what time you come to do my rooms, I will try to be here and tell you a bit of it every day until it is finished. It's a lovely long one—and I'm always putting new bits to it.«

»Then,« breathed Becky, devoutly, »I wouldn't mind \textsc{how} heavy the coal boxes was—or \textsc{what} the cook done to me, if—if I might have that to think of.«

»You may,« said Sara. »I'll tell it \textsc{all} to you.«

When Becky went downstairs, she was not the same Becky who had staggered up, loaded down by the weight of the coal scuttle. She had an extra piece of cake in her pocket, and she had been fed and warmed, but not only by cake and fire. Something else had warmed and fed her, and the something else was Sara.

When she was gone Sara sat on her favourite perch on the end of her table. Her feet were on a chair, her elbows on her knees, and her chin in her hands.

»If I \textsc{was} a princess—a \textsc{real} princess,« she murmured, »I could scatter largess to the populace. But even if I am only a pretend princess, I can invent little things to do for people. Things like this. She was just as happy as if it was largess. I'll pretend that to do things people like is scattering largess. I've scattered largess.«

